\chapter{Rekonstruksi Kurikulum}

Bab ini memuat perbandingan struktur kurikulum MBKM 2018, 2019, 2020, dan Kurikulum OBE 2025 sebagai dokumentasi jejak perubahan kurikulum.

\section{Ringkasan Rekonstruksi 2018, 2019, dan 2020}

\begin{tabular}{L{3cm}L{10cm}}
\toprule
Tahun & Catatan Rekonstruksi \\
\midrule
2018 & Penambahan skema magang, pengembangan karir, dan penyesuaian struktur proyek MBKM. \\
2019 & Penyesuaian nama mata kuliah, penambahan QMS, komputasi awan, serta aktivitas MBKM semester 7 dan 8. \\
2020 & Penguatan critical thinking, kewirausahaan, cloud, serta aktivitas MBKM pada semester akhir. \\
2025 & Transisi ke pendekatan \textit{Outcome-Based Education} (OBE) mengikuti Nota Dinas Wakil Direktur I No. 19/WADIR.I/KM/2025 (6 Februari 2025); konsolidasi seluruh pilihan MBKM Semester 7--8 menjadi satu skema Magang (20 SKS), penambahan mata kuliah sains dasar (Fisika) dan mata kuliah spesialisasi baru (Komputasi Hijau, Administrasi dan Keamanan Jaringan, Proyek Teknologi Terintegrasi/Proyek Inovasi), serta restrukturisasi pemetaan CPL-CPMK-mata kuliah agar satu mata kuliah dapat berkontribusi pada lebih dari satu CPL. \\
\bottomrule
\end{tabular}

\section{Mata Kuliah yang Dihapus, Diubah, dan Ditambahkan}

\textit{Catatan: daftar rinci perubahan mata kuliah pada 2018--2020 berikut ditranskripsi kembali dari dokumen kurikulum 2020 (BAB VII). Beberapa kode mata kuliah dipakai ulang antar tahun rekonstruksi (mis. RTI207001 tercatat sebagai "Kapita Selekta" pada daftar dihapus dan sebagai "Magang Industri 1" pada daftar baru tahun yang sama) --- ini adalah artefak penomoran pada dokumen sumber 2020, bukan kesalahan transkripsi, dan direproduksi apa adanya.}

\subsection*{Rekonstruksi 2018}

\begin{tabular}{L{3cm}L{9cm}}
\toprule
\multicolumn{2}{l}{\textbf{a. Matakuliah Dihapus}} \\
\midrule
RTI187001 & Kapita Selekta \\
RTI187003 & Proposal Skripsi \\
RTI187004 & Sistem Terdistribusi \\
RTI187006 & Pemrograman Game \\
RTI188002 & Etika Profesi Bidang TI \\
\midrule
\multicolumn{2}{l}{\textbf{b. Matakuliah Baru}} \\
\midrule
RTI187001 & Magang Industri 1 (2 bulan) \\
RTI187002 & Magang Industri 2 (4 bulan) \\
RTI187004 & Manajemen Proyek IT \\
RTI187006 & Pengembangan Karir \\
RTI188003 & Magang Industri 3 \\
\bottomrule
\end{tabular}

\subsection*{Rekonstruksi 2019}

\begin{tabular}{L{3cm}L{9cm}}
\toprule
\multicolumn{2}{l}{\textbf{a. Matakuliah Dihapus}} \\
\midrule
RTI195001 & E-Business \\
RTI195007 & Pemrograman Jaringan \\
RTI196001 & Digital Enterpreneurship \\
RTI196004 & Komputasi Multimedia \\
RTI197001 & Kapita Selekta \\
RTI197003 & Proposal Skripsi \\
RTI197004 & Sistem Terdistribusi \\
RTI197006 & Pemrograman Game \\
RTI198002 & Etika Profesi Bidang TI \\
\midrule
\multicolumn{2}{l}{\textbf{b. Matakuliah Berubah Nama}} \\
\midrule
RTI195004 & Data Warehouse $\to$ Business Intelligence \\
RTI196003 & Teknologi Data $\to$ Big Data \\
\midrule
\multicolumn{2}{l}{\textbf{c. Matakuliah Baru}} \\
\midrule
RTI195007 & QMS \\
RTI195008 & Komputasi Awan \\
RTI197001 & Magang Industri 1 (2 bulan) \\
RTI197002 & Magang Industri 2 (4 bulan) \\
RTI197003 & KKN Tematik \\
RTI197004 & Mengajar di Sekolah \\
RTI197005 & Pertukaran Pelajar \\
RTI197006 & Penelitian \\
RTI197007 & Kegiatan Wirausaha 1 \\
RTI197008 & Proyek Independen 1 \\
RTI197009 & Proyek Kemanusiaan \\
RTI198002 & Manajemen Proyek \\
RTI198003 & Pengembangan Karir \\
RTI198004 & Magang Industri 3 \\
RTI198005 & Kegiatan Wirausaha 2 \\
RTI198006 & Proyek Independen 2 \\
\bottomrule
\end{tabular}

\subsection*{Rekonstruksi 2020}

\begin{tabular}{L{3cm}L{9cm}}
\toprule
\multicolumn{2}{l}{\textbf{a. Matakuliah Dihapus}} \\
\midrule
RTI204003 & Manajemen Proyek \\
RT20I4007 & Sistem Manajemen Basis Data \\
RTI205001 & E-Business \\
RTI205007 & Pemrograman Jaringan \\
RTI206001 & Digital Enterpreneurship \\
RTI206004 & Komputasi Multimedia \\
RTI207001 & Kapita Selekta \\
RTI207003 & Proposal Skripsi \\
RTI207004 & Sistem Terdistribusi \\
RTI207006 & Pemrograman Game \\
RTI208002 & Etika Profesi Bidang TI \\
\midrule
\multicolumn{2}{l}{\textbf{b. Matakuliah Berubah Nama}} \\
\midrule
RTI204001 & Sistem Informasi $\to$ Sistem Informasi Manajemen \\
RTI205004 & Data Warehouse $\to$ Business Intelligence \\
RTI206003 & Teknologi Data $\to$ Big Data \\
\midrule
\multicolumn{2}{l}{\textbf{c. Matakuliah Baru}} \\
\midrule
RTI203006 & Matematika 3 \\
RTI203009 & QMS \\
RTI204001 & Critical Thinking dan Problem Solving \\
RTI205001 & Kewirausahaan \\
RTI205008 & Komputasi Awan \\
RTI207001 & Magang Industri 1 (2 bulan) \\
RTI207002 & Magang Industri 2 (4 bulan) \\
RTI207003 & KKN Tematik \\
RTI207004 & Mengajar di Sekolah \\
RTI207005 & Pertukaran Pelajar \\
RTI207006 & Penelitian \\
RTI207007 & Kegiatan Wirausaha 1 \\
RTI207008 & Proyek Independen 1 \\
RTI207009 & Proyek Kemanusiaan \\
RTI208002 & Manajemen Proyek \\
RTI208003 & Pengembangan Karir \\
RTI208004 & Magang Industri 3 \\
RTI208005 & Kegiatan Wirausaha 2 \\
RTI208006 & Proyek Independen 2 \\
\bottomrule
\end{tabular}

\subsection*{Rekonstruksi 2025}

Kurikulum 2025 Program Studi D4 Teknik Informatika dirancang dengan menerapkan prinsip \textit{Outcome-Based Education} (OBE) dan \textit{student-centered learning}, selaras dengan standar akreditasi IABEE (\textit{Indonesian Accreditation Board for Engineering Education}). Penyelarasan ini tercermin pada penyusunan Capaian Pembelajaran Lulusan (CPL) dan Capaian Pembelajaran Mata Kuliah (CPMK) pada Bab IV, serta pada Rencana Pembelajaran Semester (RPS) seluruh mata kuliah pada Lampiran II yang telah disusun berbasis OBE.

\begin{tabular}{L{3cm}L{9cm}}
\toprule
\multicolumn{2}{l}{\textbf{a. Matakuliah Dihapus}} \\
\midrule
RTI201002 & Teknik Dokumentasi \\
RTI201003 & Ilmu Komunikasi dan Organisasi \\
RTI201004 & Aplikasi Komputer Perkantoran \\
RTI201007 & Matematika Diskrit \\
RTI203006 & Matematika 3 \\
RTI203009 & QMS \\
RTI205002 & Proyek 2 \\
RTI207003 & KKN Tematik \\
RTI207004 & Mengajar di Sekolah \\
RTI207005 & Pertukaran Pelajar \\
RTI207006 & Penelitian \\
RTI207007 & Kegiatan Wirausaha 1 \\
RTI207008 & Proyek Independen 1 \\
RTI207009 & Proyek Kemanusiaan \\
RTI208005 & Kegiatan Wirausaha 2 \\
RTI208006 & Proyek Independen 2 \\
RTI207001 & Magang Industri 1 (2 bulan)$^{\dagger}$ \\
RTI207002 & Magang Industri 2 (4 bulan)$^{\dagger}$ \\
RTI208004 & Magang Industri 3$^{\dagger}$ \\
\midrule
\multicolumn{2}{l}{\textbf{b. Matakuliah Berubah Nama}} \\
\midrule
RTI203003 & Kecerdasan Buatan $\to$ Kecerdasan Artifisial \\
RTI204004 & Machine Learning $\to$ Pembelajaran Mesin \\
RTI205008 & Komputasi Awan $\to$ Cloud Computing \\
RTI205001 & Kewirausahaan $\to$ Kewirausahaan Berbasis Teknologi \\
RTI205005 & Pengujian Perangkat Lunak $\to$ Penjaminan Mutu Perangkat Lunak \\
RTI204002 & Proyek 1 $\to$ Proyek Sistem Informasi \\
\midrule
\multicolumn{2}{l}{\textbf{c. Matakuliah Baru}} \\
\midrule
RTI251009 & Fisika \\
RTI253006 & Metode Numerik \\
RTI255008 & Administrasi dan Keamanan Jaringan \\
RTI256001 & Komunikasi dan Etika Profesi \\
RTI256104 & Proyek Teknologi Terintegrasi \\
RTI256108 & Komputasi Hijau \\
RTI256204 & Proyek Inovasi \\
RTI256205 & Workshop Teknologi Terapan \\
RTI256206 & Rekayasa Sistem \\
RTI256207 & Teknologi Terapan \\
RTI257001 & Magang$^{\dagger}$ \\
\bottomrule
\end{tabular}

\textit{Catatan: tabel di atas disusun dari perbandingan Struktur Kurikulum MBKM 2020 (dokumen kurikulum 2020 BAB VII, "7.3 Struktur Kurikulum MBKM 2020") dengan Kurikulum OBE Rekonstruksi 2025. Hanya mata kuliah yang benar-benar berubah (dihapus, berganti nama, atau baru) yang dicantumkan, mengikuti konvensi tabel tahun-tahun sebelumnya; mata kuliah yang berlanjut tanpa perubahan berarti tidak didaftar ulang. Tanda $^{\dagger}$ menandai konsolidasi tiga skema magang bertahap 2020 (Magang Industri 1--3, total 3+ bulan berjenjang) menjadi satu mata kuliah Magang 20 SKS pada 2025 --- bukan penghapusan aktivitas magang itu sendiri, melainkan perubahan struktur penjadwalannya.}

Rekonstruksi 2025 dilaksanakan berdasarkan Nota Dinas Wakil Direktur I No. 19/WADIR.I/KM/2025 tanggal 6 Februari 2025 sebagai dasar formal transisi ke pendekatan OBE, yang kemudian diterjemahkan menjadi penyusunan ulang CPL, CPMK, dan struktur mata kuliah sebagaimana didokumentasikan pada bab-bab sebelumnya.
